\hypertarget{r29-an1-an1-forward-hnm-common-bulk-boundary-comparison}{%
\section{HNM AN1: common-bulk boundary comparison\workstar}\label{r29:an1:r29-an1-an1-forward-hnm-common-bulk-boundary-comparison}}
% BEGIN CURRENT REGISTRY NOTE
\paragraph{Statement alias HNM-P-AN1.}\label{stmt:AN1} This identifier denotes the scoped mathematical claim admitted by this chapter\textquotesingle s gate. It adds no assumption and establishes no literature-priority claim.
\begin{quote}\small\textbf{Current extension: HNM-C-AN2.} The historical limitation below remains the scope of this original result. The later, separately gated statement in Section~\ref{stmt:AN2} records the extension: With the original orthant smallness and independent HTW condition retained, deep-bulk translates of the actual orthant state converge in local state norm to a compatible locally normal centered-box state on every fixed finite complete-factor region F. The uniform enclosure is min(2,exp(C1|F|-C2(n-r-1))) for F contained in [-r,r]\textasciicircum{}3. Independent outer cutoffs, all-pairs centered Cauchy convergence, compatibility, full gauge invariance and coarse translation invariance are established before identifying the ordered limits. Its own model and state assumptions remain required.\end{quote}
% END CURRENT REGISTRY NOTE

\begin{quote}\small \textbf{Admission: accepted within scope.} Under the separately retained HTW smallness condition, the translated actual orthant box \([-n,n]^{3}\) and centered box \([-2n,2n]^{3}\) obey the uniform local bound min(2,exp(2C1-C2(n-2))) on the unit ball of \(B(H_Y)\), \(Y=\{0,e_z\}\), \(n\ge 3\). The proof retains the actual whole-star boundary, all 48 links and 36 endpoints, and the domain-preserving bounded-commutator ground condition.\par The typeset body below preserves the source-bound forward derivation. Its prospective language is historical; this boxed gate records the final reviewed verdict. The separately authored reverse and skeptical records remain linked. The common admission and attributed refinements are determined by the gate, not by a renamed equation.\end{quote}
Human project author: \textbf{Hruday N M (BUNZEEY)}. Independent forward execution under the frozen contract, without current reverse or skeptic results. The theorem below applies an existing locality theorem to this exact boundary dictionary; scientific priority is unverified.

\textbf{Finite comparison:} write \(A_n=[-n,n]^{3}\), the coarse translate of \([0,2n]^{3}\), and \(B_n=[-2n,2n]^{3}\). For the contracted identical coarse-homogeneous selected coefficient triple and omitted tau, the actual whole-star ground states obey, for \(Y=\{0,e_z\}\) and \(n\ge 3\),

\begin{equation}
\sup_{\|A\|\le1,\ A\in B(\mathcal H_Y)}
|\rho_{A_n}(A)-\rho_{B_n}(A)|
\le\min\{2,\exp(2C_1-C_2(n-2))\}.
\tag{HNM-AN1.1}
\label{r29:an1:eq1}
\end{equation}

This uses the \textbf{additional} source smallness \(7|\tau|\le c_{\mathrm{HTW}}(1,1)\), with positive unevaluated source constants. They are not AM2's numerical radius or I1's c1,c2. Infinite-state identification is not this loop's result.

\hypertarget{r29-an1-common-hilbert-space-and-canonical-interior}{%
\subsection{Common Hilbert space and canonical interior}\label{r29:an1:r29-an1-common-hilbert-space-and-canonical-interior}}

Extend \(\rho_{A_n}\) to \(B_n\) by the actual onsite vacuum product outside \(A_n\); call the resulting normal state rho1. It is the ground state of

\(H_1=H_{A_n}+\sum_{x\in B_n\setminus A_n}h_x\).

Let rho2 be the actual whole-star ground of \(H_2=H_{B_n}\). Both are trace-one densities on the \textbf{same} finite tensor Hilbert space. The original compact elliptic rotor ground exists and is unique/positive for the bounded real finite interaction; this assertion does not require guessing a theorem constant.

Set \(\Lambda_*=A_n\) and define the source canonical restriction

\begin{equation}
H_*=\sum_{x\in A_n}h_x+
\sum_{x:\operatorname{dist}_1(x,A_n^c)>1}\phi_x.
\tag{HNM-AN1.2}
\label{r29:an1:eq2}
\end{equation}

The retained source anchors are exactly \([-n+1,n-1]^{3}\). Each source \(\phi_x\) is supported in the radius-one l1 ball about x because the whole star S is contained in that ball, even though S has l1 diameter two. Its norm is \(M=7|\tau|\), and the onsite normalized gap is at least one. This matches HTW Definition1 with \(R=1\) and \(g=1\); \(R=0\) would discard actual neighboring factors.

The source bulk from Definition6 is the twice-eroded set

\begin{equation}
D_n=A_n^\circ=\{x:\operatorname{dist}_1(x,A_n^c)>2\}
=[-n+2,n-2]^3.
\tag{HNM-AN1.3}
\label{r29:an1:eq3}
\end{equation}

Actual whole-star anchors of \(A_n\) are \([-n,n-1]^{3}\), not the canonical source anchors. Their extra groups have some base coordinate -n, and every point of their star has that coordinate at most -n+1; therefore their support misses D\_n.~For H2, every anchor not in the canonical set has a coordinate \(\le -n\) or \textgreater=n.~Its star has that coordinate \(\le -n+1\) or \(\ge n\), again outside D\_n.~Onsite terms outside \(A_n\) also miss D\_n.~Thus each \(H_i-H_*\) is supported outside \(D_n\), including all missing/extra incoming stars. Equality of boundary Hamiltonians has not been assumed.

\hypertarget{r29-an1-unbounded-domain-ground-condition}{%
\subsection{Unbounded-domain ground condition}\label{r29:an1:r29-an1-unbounded-domain-ground-condition}}

It remains to prove the source's ground-in-the-bulk hypothesis, rather than equate states from a common gap. A source test operator T is bounded and supported in \(D_n\), preserves \(D(H_\{0,D_n\})\), and has bounded commutator with that local onsite sum. For a finite sum of nonnegative commuting tensor onsite operators,

\(D(H_{0,B_n})=D(H_{0,D_n})\cap D(H_{0,B_n\setminus D_n})\).

Identity extension of T commutes with the exterior spectral projections and preserves the latter domain. It consequently preserves the full domain. Every finite interaction is bounded, so each \(H_i\) has that same onsite domain. Every interaction in \(H_i-H_*\) acts in the disjoint tensor factor outside \(D_n\), and commutes with T; the exterior onsite terms commute on the stated domain as well. Hence

\begin{equation}
[H_i,T]=[H_*,T]\quad\hbox{as bounded extensions},\qquad i=1,2.
\tag{HNM-AN1.4}
\label{r29:an1:eq4}
\end{equation}

Their genuine ground vectors Psi\_i therefore give

\(\Tr(\rho_iT^*[H_*,T])=\langle T\Psi_i,(H_i-E_i)T\Psi_i\rangle\ge0\).

This is exactly HTW Definition6's ground-in-the-common-bulk condition. It does not assume every bounded local T preserves an unbounded generator domain. The code retains a rank-one counterexample to that false universal domain claim. The output observable A in (HNM-AN1.1), in contrast, may be any bounded operator in the local factor, as allowed by the source theorem after the ground condition has been established.

\hypertarget{r29-an1-the-theorem-and-exact-distance}{%
\subsection{The theorem and exact distance}\label{r29:an1:r29-an1-the-theorem-and-exact-distance}}

Henheik--Teufel--Wessel~\citep{henheik2022local}, arXiv:2106.13780v3, Definitions1,2,6, Theorem7 and Lemma8, permit infinite-dimensional onsite factors and a bound on two normal states ground in a common weakly interacting bulk. The required interaction threshold c\_HTW and decay constants C1,C2 are existential. Their Theorem7, attributed there to Yarotsky, applies to rho1,rho2 and \(H_*\) and gives

\(|\Tr((\rho_1-\rho_2)A)|\le e^{C_1|Y|-C_2\operatorname{dist}(Y,\mathbb Z^3\setminus D_n)}\|A\|\).

For \(Y=\{0,e_z\}\), \(|Y|=2\) and the exact l1 distance to the complement of \(D_n\) is n-2 for \(n\ge 3\) (the \(e_z\) point is nearest the upper z boundary). This proves (HNM-AN1.1). The bound is uniform over the full norm-one bounded local algebra, hence is also the trace-norm bound on the two Y-marginals and applies to its gauge-invariant subalgebra.

Y consists of \textbf{48 actual links and 36 distinct original endpoints}. Gauge actions at all those endpoints are preserved by the coarse translation \((4n,2n,n)\) in fine coordinates and by both boundary prescriptions. Coarse translation invariance is essential: the same selected triple is used in every factor. Mere membership in coefficient ranges would not identify the translated interaction.

\hypertarget{r29-an1-controls-and-limits}{%
\subsection{Controls and limits}\label{r29:an1:r29-an1-controls-and-limits}}

The checker reconstructs both boundary inventories and all star supports for \(n=3,4,5\), computes the eroded set and exact distance, and verifies that every boundary-difference group misses \(D_n. It\) reconstructs four incoming/outgoing anchors of an interior singleton and seven for Y; keeping only the two anchors in Y loses five groups. It also checks the 48-link/36-endpoint completion.

A near-boundary region with a point at \((-n,0,0)\) has distance zero to the complement of \(D_n\) for all n; this comparison gives no vanishing enclosure there. A changed interaction at a bulk anchor generally fails (HNM-AN1.4); an exact matrix commutator demonstrates that support overlap cannot be dropped. Changing a translated selected coefficient also destroys the common model. The additional source smallness is retained symbolically; it cannot be deleted or replaced by the explicit AM2 gap radius without a new theorem.

The numerical fixtures are geometry and domain countercontrols, not simulated expectation values. We have proved a conditional finite comparison, and have not yet identified the actual inherited orthant state with any centered-box state. That requires a separate limit argument under the full simultaneous premises, including I1's \(|\tau|<\tau_*\) whenever that inherited state is invoked.

Reproduce from the repository root:
\begin{verbatim}
python -B research/round29/forward/an1/check.py \
  --output /absolute/new/output
\end{verbatim}

\subsection*{Final admitted limits and evidence}
\begin{itemize}
\item This loop proves a finite comparison only; diagonal finite agreement does not identify infinite-volume states.
\item The selected coefficient triple must be fixed and coarse-translation invariant. A changed bulk interaction is outside this comparison.
\item The additional HTW threshold \(7|\tau|\le c_{\mathrm{HTW}}(1,1)\) and decay constants remain unevaluated; AM2 does not supply them.
\item Invoking the actual inherited orthant state also requires its original \(|\tau|<\tau_*\) restriction. Near-boundary controls show failure of this decay certificate, not actual unequal expectations.
\item This is a model-specific application of HTW/Yarotsky, not a continuum construction or an established priority claim.
\end{itemize}
\repo{research/round29/advisor/an1-gate.json}\par\repo{research/round29/reverse/an1/report.md}\par\repo{research/round29/skeptic/an1.md}

\emph{Sequencing note.} Prospective statements in this frozen derivation describe the moment of its admission. The final Round29 roadmap below governs subsequent work.
